% EPPA-III Section 6 / Proposition 6.17.
% An optional minimal source-local bridge; NOT an edit to the frozen manuscript.
% Mathematical conclusion separately checked by the stronger Lean theorem
% EPPAIII.CliqueCovers.eppa_noncluster_clique_double_all_positive, PR #307.
% The paragraph below supplies the missing outside-half hypothesis of Lemma 6.1
% without changing the proof of that lemma. Author review pending.

\paragraph{The complementary half in the equal-clique case.}
Suppose that $H$ is non-homogeneous and vertex-transitive,
$|H|=2st$, and $G\cong sK_t$ is an induced subgraph for which
$H$ is an EPPA witness. Put $B=H\setminus G$ and denote the
cliques of $G$ by $C_1,\ldots,C_s$. We claim that
$B\cong sK_t$, so that Lemma~\ref{thm:sKt_witnesses} applies.

Since $H$ is regular and the two halves have the same order,
counting their internal and crossing edges gives
$|E(B)|=|E(G)|=s\binom t2$.
The cases $s=1$ and $t=1$ follow immediately.
Suppose $s,t\ge2$. For $w\in B$, let
\[
  P(w)=\bigl(N_H(w)\cap C_i\bigr)_{i=1}^{s}.
\]
The extensions of $\Aut(G)=S_t\wr S_s$ show that the orbit
of every such profile is represented in $B$, and hence has
size at most $st$. If $p$ coordinates are partial, the orbit
has size at least $\binom sp t^p$. This exceeds $st$ for
$p\ge2$, except when $s=t=p=2$. When $p=1$, with $k$
neighbours in the partial coordinate and $r$ full coordinates,
the exact orbit size is
$s\binom tk\binom{s-1}r\le st$. Hence
$k\in\{1,t-1\}$ and $r\in\{0,s-1\}$.
In both cases the orbit has precisely $st$ elements,
and gives canonical labels of all vertices of $B$.
The induced action on $B$ has two orbits on unordered pairs
(the two pairs of opposite corners in the exceptional square
case). The equality of internal edge counts therefore gives
$B\cong sK_t$.

It remains to assume that every profile has only empty
or full coordinates. Then each $C_i$ consists of true twins
of $H$. By vertex-transitivity the true-twin classes all
have the same order $a\ge t$. The quotient on these classes
has at most $2s$ vertices and contains an independent
set of $s$ vertices. It has positive degree, since a
vertex-transitive disjoint union of cliques is homogeneous.
A regular graph of positive degree has no independent
set of more than half of its vertices. Consequently
$a=t$ and the quotient has exactly $2s$ vertices.
The $s$ remaining classes constitute $B$, and the
internal edge count forces them to be pairwise
nonadjacent. Again $B\cong sK_t$, as claimed.

% Rationale: Proposition 6.17 currently invokes Lemma 6.1 before
% establishing its complementary-half hypothesis. This local claim
% supplies precisely that missing premise while preserving the lemma.
% It uses no Gardiner or rank-three classification, and makes no
% choice of coherent EPPA extensions. The complete profile argument
% is verified at theorem level in PR #307; this compact paper
% argument must still be reviewed against the source's local notation.
